\documentclass[10pt,twocolumn]{article}

\usepackage[utf8]{inputenc}
\usepackage[T1]{fontenc}
\usepackage{times}
\usepackage{graphicx}
\usepackage{amsmath}
\usepackage{amssymb}
\usepackage{booktabs}
\usepackage{hyperref}
\usepackage{subcaption}
\usepackage[accsupp]{axessibility}

\graphicspath{{../figures/}}
\newcommand{\argmax}{\operatornamewithlimits{argmax}}

\title{Memory Update Decisions as Constrained Resource Allocation in 3D Reconstruction}

\author{Agostina Silva\\
Independent Researcher\\
{\tt\small agostina.dev.dc@gmail.com}
}

\date{}

\begin{document}

\maketitle

\begin{abstract}
Real-time 3D reconstruction systems (e.g.\ 3D Gaussian Splatting) face a fundamental resource allocation problem: given thousands to millions of scene entities and limited per-frame compute, which entities should be updated? We show that entity update benefits are \textit{non-separable}---the value of updating entity A depends on whether B is also updated---violating the independence assumption underlying greedy selection. We formalize this as a budgeted subset selection problem, measure coupling ($\kappa$) across entity pairs in 6 real 3DGS scenes (bicycle, bonsai, counter, garden, room, kitchen), and train a lightweight Q-network that learns to predict marginal update benefit. Mean $|\kappa| = 0.0008$ across scenes (range $0.0006$--$0.0012$), with $22.3\%$ of pairs coupled on average. The learned Q-function outperforms a motion/uncertainty heuristic by $+17.1\%$ on average. A CEM planner operating on top of Q-values further improves this to $+22.2\%$. These results demonstrate that non-separable coupling is a general phenomenon in 3DGS and can be exploited for better budgeted entity selection.
\end{abstract}

\noindent\textbf{Code:} \url{https://github.com/Agostynah/memory-update-3dgs}

\section{Introduction}
\label{sec:intro}

3D Gaussian Splatting (3DGS)~\cite{kerbl3dgs} represents scenes as collections of anisotropic 3D Gaussians. Real-time rendering iterates over all Gaussians each frame. Under computational budgets typical of mobile devices, robotics, or VR headsets, updating every entity every frame is infeasible. The system must choose a subset of $K$ entities to update, given a per-frame budget.

Current practice relies on heuristics: update entities with largest motion, highest uncertainty, or simply random selection. These approaches implicitly assume that entity selection is \textit{separable}---the benefit of updating entity $i$ does not depend on which other entities are also updated. We demonstrate that this assumption is false.

\textbf{Key insight.} When multiple 3D Gaussians overlap, their contributions to the rendered image are superposed. Updating entity $A$ changes the same pixels that $B$ affects. The combined benefit of updating both is not the sum of individual benefits. This \textit{coupling} means greedy selection by individual scores can be arbitrarily suboptimal~\cite{krause2005near}.

We make the following contributions:
\begin{itemize}
    \item We formalize memory update in 3D reconstruction as a budgeted subset selection problem and define coupling $\kappa$ to measure non-separability.
    \item We demonstrate $\kappa > 0$ across 6 real 3DGS scenes (mean $|\kappa|=0.0008$, $22.3\%$ of entity pairs coupled).
    \item We train a lightweight Q-network (3-layer MLP, $\sim$40KB) that predicts per-entity update benefit and selects entities at inference.
    \item The learned Q-function outperforms a strong heuristic by $+17.1\%$; a CEM planner on top of Q-values improves this to $+22.2\%$.
    \item Ablation study confirms that cluster identity is the most critical feature for Q-network accuracy.
\end{itemize}

\section{Related Work}

\textbf{3DGS efficiency.} Several works address the computational cost of 3DGS. Lee et al.~\cite{lee2024compact} propose compact scene representation via pruning and quantization. Yu et al.~\cite{yu2024mcmc} use MCMC sampling for densification control. Shen et al.~\cite{shen2025lodgs} introduce LOD-based rendering for large scenes. These methods optimize the static scene representation but do not address per-frame update selection under budget constraints.

\textbf{Budgeted subset selection.} The problem of selecting a subset of items subject to cardinality constraints is well-studied in submodular optimization~\cite{krause2005near,wei2015constrained}. However, submodularity---the property of diminishing returns that enables efficient greedy algorithms---is assumed, not verified. We show that entity update benefits are non-monotone and potentially non-submodular, motivating learned selection.

\textbf{Learning for resource allocation.} Reinforcement learning has been applied to sensor scheduling~\cite{leong2020deep}, caching~\cite{sadeghi2019deep}, and neural network layer selection~\cite{wu2018blockdrop}. To our knowledge, ours is the first application to entity update decisions in 3D reconstruction.

\section{Problem Formulation}

Consider a scene with $N$ entities (Gaussians) tracked over $T$ frames. At each frame $t$, we observe noisy entity states $G_t^{\text{obs}}$ (rendered from noisy parameters) and have access to ground truth $G_t^*$ (rendered from true parameters). We have budget $K \ll N$ updates.

Define the benefit of updating a subset $S \subseteq \{1,\ldots,N\}$:

\begin{equation}
R(S, t) = \text{MSE}(G_t^*, G_t^{\text{obs}}) - \text{MSE}(G_t^*, G_t^{(S)})
\end{equation}

where $G_t^{(S)}$ is the rendering after updating exactly entities in $S$ to their true parameters.

\textbf{Non-separability.} For entities $i$ and $j$, define pairwise coupling:

\begin{equation}
\kappa_{ij}(t) = R(\{i,j\}, t) - (R(\{i\}, t) + R(\{j\}, t))
\end{equation}

$\kappa_{ij} = 0$ indicates separable benefits; $\kappa_{ij} \neq 0$ indicates coupling. When $\kappa_{ij} > 0$, the joint benefit exceeds the sum of individual benefits (complementarity). When $\kappa_{ij} < 0$, there is redundancy.

\textbf{Selection objective.} At each frame $t$, given budget $K$, solve:

\begin{equation}
S_t^* = \argmax_{S: |S| = K} R(S, t)
\end{equation}

We compare three approaches:
\begin{itemize}
    \item \textbf{Oracle:} brute-force evaluation of $\binom{N}{K}$ subsets (feasible for $K \leq 4$).
    \item \textbf{Greedy heuristic:} rank by $h(i,t) = \|\Delta G_i\| + \|G_i^{\text{obs}} - G_i^*\| + (1 - c_i)$ where $\Delta G_i$ is motion and $c_i$ is confidence.
    \item \textbf{Learned Q:} rank by $Q_\theta(i,t)$, a learned predictor of $R(\{i\}, t)$.
\end{itemize}

\section{Method}

\subsection{Voxel Proxy World}

We construct a controlled 3D environment. The world is a $16^3$ voxel grid. Each entity is a 3D Gaussian with position $\mu_e \in \mathbb{R}^3$, isotropic scale $\sigma_e$, and amplitude $a_e$. The rendered value at voxel $(x,y,z)$ is:

\begin{equation}
v(x,y,z) = \sum_{e=1}^N a_e \exp\left(-\frac{\|(x,y,z) - \mu_e\|^2}{2\sigma_e^2}\right)
\end{equation}

$\text{MSE}(a,b) = \frac{1}{G^3} \sum_{x,y,z} (a_{xyz} - b_{xyz})^2$.

Entities move via Brownian motion with cluster attraction (6 clusters). Observations add Gaussian noise scaled per entity. $N=80$, $T=120$, yielding 9600 training samples.

\textbf{Real data extension.} We download 6 pretrained 3DGS scenes from the Mip-NeRF 360 dataset~\cite{mipnerf360} (bicycle, bonsai, counter, garden, room, kitchen; 273K--4.4M Gaussians each, obtained from \texttt{dylanebert/3dgs}). For each scene, filter valid Gaussians and run KMeans ($k=80$) to form meta-entities preserving spatial structure. Meta-entity position is the cluster centroid; size is proportional to cluster radius. Motion and observation noise follow the same generative process as the synthetic world, allowing consistent evaluation across scenes.

\begin{table}[t]
\centering
\caption{Entity parameters across 6 real 3DGS scenes. Scenes sorted by coupling strength.}
\label{tab:params}
\begin{tabular}{lrrr}
\toprule
Scene & Gaussians (K) & $|\kappa|$ mean & Coupled \% \\
\midrule
room & 1130 & 0.0006 & 16.5\% \\
kitchen & 1685 & 0.0006 & 17.6\% \\
bonsai & 273 & 0.0007 & 17.6\% \\
counter & 1029 & 0.0008 & 22.2\% \\
bicycle & 568 & 0.0009 & 27.1\% \\
garden & 4386 & 0.0012 & 32.6\% \\
\midrule
Mean & -- & 0.0008 & 22.3\% \\
\bottomrule
\end{tabular}
\end{table}

\subsection{Feature Design}

For each entity $e$ at frame $t$, we compute 16 features:
\begin{itemize}
    \item \textbf{Position (3):} $\mu_e / G$ (normalized to $[0,1]$)
    \item \textbf{Motion (3):} $\mu_e(t) - \mu_e(t-1)$ if $t>0$, else zero
    \item \textbf{Uncertainty (1):} $\|G_t^{\text{obs}} - G_t^*\| / G$
    \item \textbf{Confidence (1):} $c_e = 1 - 0.3 \cdot d_e / \sigma_e$ where $d_e$ is distance from cluster center
    \item \textbf{Scale (1):} $\sigma_e / G$
    \item \textbf{Amplitude (1):} $a_e$
    \item \textbf{Overlap count (1):} fraction of entities with pairwise IoU $> 0.05$
    \item \textbf{Cluster ID (5):} one-hot encoding of cluster membership
\end{itemize}

\subsection{Q-Network}

We train a small MLP to predict $R(\{e\}, t)$:

\[
16 \to 64 \to 32 \to 1,
\quad \text{ReLU, Dropout(0.1)}
\]

Optimizer: AdamW (lr=1e-3, weight decay=0.01). Targets standardized to zero mean, unit variance. Early stopping with 200 epochs max. Training takes $\sim$5 minutes on CPU. Inference is $\sim$0.1ms per entity on CPU.

\subsection{CEM Planner}
\label{sec:cem}

Independent Q-scoring cannot capture pairwise coupling during selection. To address this, we augment Q-scores with a Cross-Entropy Method (CEM) planner~\cite{rubinstein1999cross}. CEM iteratively refines a sampling distribution over entity subsets:

\begin{enumerate}
    \item Initialize a categorical proposal distribution $\mathbf{p}^{(0)}$ proportional to $\max(Q_\theta(i), 0) + \epsilon$.
    \item For iteration $m = 1,\ldots,M$:
    \begin{enumerate}
        \item Draw $S$ candidate subsets of size $K$ from $\mathbf{p}^{(m-1)}$.
        \item Evaluate each subset via the world model $R(S, t)$.
        \item Set $\mathbf{p}^{(m)}$ to the empirical distribution of the top $E$ elite subsets (weighted combination: $0.7\mathbf{p}^{(m)} + 0.3\mathbf{p}^{(m-1)}$).
    \end{enumerate}
    \item Return the best subset found.
\end{enumerate}

We use $M=4$ iterations, $S=80$ samples, $E=16$ elites. CEM adds $\sim$50ms per budget level, negligible compared to per-frame rendering cost.

\section{Experiments}

\paragraph{Coupling measurement ($\kappa$).}
For each dataset, we compute $\kappa_{ij}$ for all pairs at $t=0$. We report mean $|\kappa|$ and fraction of pairs exceeding $|\kappa| > 0.001$.

\paragraph{Regret.}
For $K = 2,3,4,5$, we compare greedy selection against brute-force oracle over 20 random frames (sampling 3000 subsets when $\binom{N}{K} > 3000$).

\paragraph{Q-learning.}
Train Q-network on 9600 samples (80 entities $\times$ 120 frames). Evaluate on held-out 30 frames.

\paragraph{Selection quality.}
At each budget $K=1,\ldots,8$, select top-K entities by heuristic score, Q-score, and oracle (for $K \leq 4$). Report average MSE across 30 evaluation frames.

\section{Results}

\begin{table}[t]
\centering
\caption{Main results across 6 real 3DGS scenes. Q and CEM improvement measured against the motion/uncertainty heuristic. Oracle gap computed at $K=4$ using 2000 sampled subsets.}
\label{tab:results}
\begin{tabular}{lcccccc}
\toprule
Dataset & $|\kappa|$ mean & Val loss & Q improv. & CEM improv. & Oracle gap (K=4) & $Q > H$ \\
\midrule
bicycle & 0.0009 & 0.37 & $+17.0\%$ & $+22.6\%$ & $-6.0\%$ & All K \\
bonsai & 0.0007 & 0.44 & $+18.5\%$ & $+23.7\%$ & $-3.8\%$ & All K \\
counter & 0.0008 & 0.36 & $+13.3\%$ & $+17.1\%$ & $-5.5\%$ & All K \\
garden & 0.0012 & 0.38 & $+14.1\%$ & $+19.7\%$ & $-4.6\%$ & All K \\
room & 0.0006 & 0.43 & $+21.5\%$ & $+25.6\%$ & $-7.3\%$ & All K \\
kitchen & 0.0006 & 0.37 & $+18.4\%$ & $+24.4\%$ & $-4.3\%$ & All K \\
\midrule
Mean & 0.0008 & 0.39 & $+17.1\%$ & $+22.2\%$ & $-5.2\%$ & 6/6 \\
$\pm$Std & 0.0002 & 0.04 & $\pm 2.8\%$ & $\pm 2.9\%$ & -- & -- \\
\bottomrule
\end{tabular}
\end{table}

\subsection{Coupling Across Scenes}

\begin{figure}[t]
\centering
\includegraphics[width=0.7\textwidth]{fig_coupling_overview.png}
\caption{Coupling strength ($|\kappa|$ mean) and fraction of coupled entity pairs across 6 real 3DGS scenes.}
\label{fig:coupling}
\end{figure}

All 6 real 3DGS scenes exhibit measurable coupling. Mean $|\kappa| = 0.0008$ (range $0.0006$--$0.0012$). Fig.~\ref{fig:coupling} and Table~\ref{tab:params} show the per-scene breakdown. The fraction of coupled entity pairs ranges from $16.5\%$ (room) to $32.6\%$ (garden), with a scene average of $22.3\%$.

The scene with the strongest coupling is garden ($|\kappa|=0.0012$, $32.6\%$ of pairs), which also has the most Gaussians (4.4M). This suggests coupling strength may increase with scene complexity, as denser Gaussians produce more spatial overlap.

\subsection{Q-Learning Captures Structure}

Validation loss averages 0.39 (range 0.36--0.44) against a baseline of 1.0 (predicting mean). Performance is consistent across scenes, with no scene worse than $-56\%$ of baseline. The Q-network generalizes across diverse scene geometries (indoor rooms, outdoor gardens, countertop objects).

\subsection{Learned Selection Beats Heuristic}

\begin{figure}[t]
\centering
\includegraphics[width=0.85\textwidth]{fig_per_scene_mse.png}
\caption{Per-scene selection quality. Each panel shows MSE vs budget $K$ for heuristic, Q-Learned, CEM, and a sampled oracle baseline.}
\label{fig:per_scene}
\end{figure}

Q-based selection outperforms the motion+uncertainty heuristic at every budget $K=1..8$ in all 6 scenes. Fig.~\ref{fig:per_scene} shows the per-scene curves; Fig.~\ref{fig:aggregated} shows the cross-scene average. Average improvement: $+17.1\% \pm 2.8\%$. The gain is monotonic in $K$, rising from $+4.9\%$ at $K=1$ to $+27.4\%$ at $K=8$. This reveals an important property: \textit{better scoring compounds}---correct ordering at higher budgets matters more because more entities are affected.

\begin{figure}[t]
\centering
\includegraphics[width=0.7\textwidth]{fig_aggregated.png}
\caption{Aggregated selection quality averaged across 6 scenes. Shaded regions indicate $\pm 1$ std.}
\label{fig:aggregated}
\end{figure}

\subsection{CEM Improves Over Independent Q-Scoring}

The CEM planner adds $+5.1\%$ on average over Q-scoring ($+22.2\%$ vs $+17.1\%$). Fig.~\ref{fig:cem_results} shows per-scene breakdowns and per-budget trends.

\begin{figure}[t]
\centering
\begin{subfigure}{0.49\textwidth}
\includegraphics[width=\textwidth]{fig_improvement.png}
\caption{Improvement by scene.}
\label{fig:improvement}
\end{subfigure}
\hfill
\begin{subfigure}{0.49\textwidth}
\includegraphics[width=\textwidth]{fig_per_k_improvement.png}
\caption{Improvement by budget $K$.}
\label{fig:per_k}
\end{subfigure}
\caption{CEM consistently outperforms independent Q-scoring. The Q contribution grows with $K$; CEM adds further gains at every budget.}
\label{fig:cem_results}
\end{figure} The improvement is largest at intermediate budgets ($K=3$--$5$), where pairwise coupling is most relevant (enough entities to interact, not so many that selection is trivial). At $K=1$, CEM reverts to Q-ranking (no coupling possible). At $K=8$, CEM still improves ($+34.5\%$ vs $+27.4\%$), suggesting coupling affects selection at all budgets.

\subsection{Oracle Gap}

The gap between Q-based selection and the sampled oracle (2000 draws from $\binom{80}{4}=1.58\text{M}$) averages $-5.2\%$ at $K=4$. This negative gap---Q appearing \textit{better} than oracle---indicates that the brute-force sampling is too sparse to reliably find the optimum. CEM achieves larger negative gaps ($-11.3\%$), confirming that both methods find better subsets than 2000 random draws. This is a measurement limitation, not a structural one.

\subsection{Ablation Study}

\begin{figure}[t]
\centering
\begin{subfigure}{0.49\textwidth}
\includegraphics[width=\textwidth]{fig_ablation.png}
\caption{Ablation on bicycle scene.}
\label{fig:ablation}
\end{subfigure}
\hfill
\begin{subfigure}{0.49\textwidth}
\includegraphics[width=\textwidth]{fig_varying_n.png}
\caption{Effect of entity count $N$.}
\label{fig:varying_n}
\end{subfigure}
\caption{Feature ablation and sensitivity to entity count. Cluster ID is the most important feature; $N=80$ is optimal.}
\label{fig:ablation_results}
\end{figure}

We ablate the Q-network on the bicycle scene (Fig.~\ref{fig:ablation_results}a). The cluster ID feature is most important ($+0.048$ val loss increase when removed), followed by dropout regularization ($+0.006$), and overlap information ($+0.003$). Varying entity count $N$ (Fig.~\ref{fig:ablation_results}b) shows that $N=80$ achieves the best val loss ($0.360$), while $N=40$ is substantially worse ($0.671$) and $N=120$ slightly degrades ($0.387$), suggesting an optimal entity count for the proxy world.

\section{Discussion}

\paragraph{Limitations.} Our voxel proxy world captures the coupling mechanism through additive Gaussian overlaps in a shared grid, but does not replicate full 3DGS rendering (anisotropic Gaussians, spherical harmonics, CUDA rasterization). The real 3DGS experiment uses meta-entities (KMeans clusters) rather than individual Gaussians, which may obscure per-Gaussian coupling patterns. The regret measurement is limited by combinatorial sampling at higher budgets.

\paragraph{Generalizability.} The coupling mechanism we identify---overlapping entities sharing the same output space---is generic. We verified $\kappa > 0$ across 6 diverse real 3DGS scenes spanning indoor, outdoor, and object-centric capture. Any 3D representation with spatial overlap (point clouds with splatting, NeRF with ray marching) likely exhibits analogous non-separability. Our Q-learning approach transfers across scenes with consistent improvement ($+17.1\% \pm 2.8\%$).

\paragraph{Future work.} Scaling to per-Gaussian decisions ($N > 10^5$) requires sparse $\kappa$ computation (e.g., spatial proximity thresholds). The CEM planner could be replaced with a learned subset decoder for real-time inference without iterative sampling. Extending to sequential decisions across frames (full RL rather than per-frame scoring) is a natural next step.

\section{Conclusion}

We have shown that entity update decisions in 3D reconstruction exhibit non-separable coupling across 6 real 3DGS scenes (mean $|\kappa|=0.0008$, $22.3\%$ of pairs). A lightweight Q-network learns to predict marginal update benefit and achieves $+17.1\%$ improvement over a motion/uncertainty heuristic. A CEM planner further improves this to $+22.2\%$, confirming that modeling coupling during selection adds value beyond independent scoring. These results open a new direction for budgeted memory management in real-time 3D systems, with applications in robotics, VR/AR, and mobile graphics.

\bibliographystyle{plain}
\bibliography{main}

\end{document}
